\chapterhead{6}{Reporting What You Found}{Saying what you know, and how well you know it}\label{ch6}

A measurement that stays in your head helps no one. A good report lets a stranger see what you did, so they can check your work or repeat it.

\subhead{The five-line report}
Every report in this series can use the same five lines.

\begin{tcolorbox}[enhanced,colback=white,colframe=ink!60,boxrule=1pt,arc=2mm,left=8pt,right=8pt,top=5pt,bottom=5pt,before skip=7pt,after skip=7pt]
\sffamily\small
\begin{enumerate}[leftmargin=1.6em,itemsep=0.3em]
\item \textbf{What I measured.} The thing, and the quantity (length, time, mass).
\item \textbf{How I measured it.} The instrument, the unit, how it was read, how many times.
\item \textbf{My readings.} All of them, not just the one I liked best.
\item \textbf{My answer.} The average or best value, rounded sensibly, with the spread.
\item \textbf{What I kept and what it hides.} What the answer keeps, and what it leaves out.
\end{enumerate}
\end{tcolorbox}

Here is a report built from the practice numbers in Chapter 4.

\begin{worked}{A sample report (practice numbers)}
\textbf{1. What I measured.} The time a sheet of scrap paper took to fall to the floor, released flat from the height of my raised hand.\\
\textbf{2. How I measured it.} A stopwatch read to the nearest 0.1~s. One person timed ten falls with the same sheet.\\
\textbf{3. My readings (s).} 1.8, 2.1, 1.9, 2.0, 2.4, 1.9, 2.0, 2.1, 1.8, 2.0.\\
\textbf{4. My answer.} About 2.0~s. The readings ranged from 1.8~s to 2.4~s, a spread of 0.6~s.\\
\textbf{5. What I kept and what it hides.} The average and the spread are kept. The order of the readings, and which drop was which, are set aside. A different person, sheet or height may give a different answer.
\end{worked}

\subhead{Which rung of the ladder?}
Remember the ladder from Chapter 3: observed, counted, estimated, concluded. A good report keeps them apart. Look at these four sentences from the sample report. Each one sits on a different rung.

\clearpage
\begin{puzzle}{Climb the ladder}
Write O (observed), C (counted), E (estimated) or K (concluded) next to each sentence.
\begin{description}[leftmargin=1.6em,itemsep=0.3em,style=sameline]
\item[(a)] ``The sheet wobbled as it fell.'' \wline[1cm]
\item[(b)] ``I timed ten drops.'' \wline[1cm]
\item[(c)] ``A sheet like this takes about 2.0~s to fall from my reach.'' \wline[1cm]
\item[(d)] ``Paper falls slowly because air pushes up on it.'' \wline[1cm]
\end{description}
Which of these sentences did the ten readings \emph{not} test? \wline[3cm]
\end{puzzle}

\subhead{Five ways to check, five things they show}
The trust box in this book used \emph{repeat and compare}. The other series use other kinds of check. Each one can show something, and none of them shows everything.

\renewcommand{\arraystretch}{1.35}
\begin{center}
\sffamily\scriptsize
\begin{tabularx}{\textwidth}{>{\bfseries}p{2.4cm}XX}
\toprule
\textbf{Kind of check} & \textbf{What it can show} & \textbf{What it cannot show}\\
\midrule
Prove (math) & A result must hold for every case the argument covers. & That the rule fits the real world.\\
Compare with measurement (science) & A prediction matched what was measured this time. & That it will match every time and everywhere.\\
Repeat and compare (science, ecology) & A result can be repeated. & That it is true, or that the test was fair.\\
Check against how words are used (spelling and grammar) & A choice fits dictionaries and good usage. & That it is the only choice, or that a word's whole history is known.\\
Check against sources (Problem Solving) & What surviving evidence says. & That the evidence is complete.\\
\bottomrule
\end{tabularx}
\end{center}

\clearpage

\thispagestyle{fancy}\label{slop}
\subhead{Slop check}
Slop is hasty, gaudy work that forgot to take care. It can be made by a person or a machine, quickly or slowly. What matters is not who made it but whether it shows care for three things:
\begin{itemize}[leftmargin=1.6em]
\item \textbf{Salient:} Is the main point easy to see?
\item \textbf{Simple:} Is there anything that does nothing but get in the way?
\item \textbf{Relevant:} Does everything on the page help answer the question?
\end{itemize}

Here is a hasty science poster. Study it, then decide what is wrong.

\begin{center}
\begin{tikzpicture}[scale=0.97]
  \fill[magenta!35,draw=green!50!black,line width=3pt] (0,0) rectangle (11.4,7.6);
  \node[star,star points=12,fill=yellow,draw=red,line width=2pt,minimum size=2.2cm,font=\bfseries\scriptsize,rotate=15,text=blue] at (1.4,6.2) {SUPER!!!};
  \node[font=\itshape\bfseries\Large,text=cyan!60!blue,rotate=-3] at (6.4,6.8) {MeAsUrInG   is   AwEsOmE};
  \node[font=\ttfamily\tiny,text=orange] at (6.4,6.15) {(results may vary lol)};
  \foreach \x/\h/\c in {2.0/3.3/red,3.0/1.1/green,4.0/4.4/blue,5.0/2.2/orange,6.0/3.9/violet,7.0/0.8/teal,8.0/2.9/pink,9.0/4.0/brown}{
    \fill[\c,draw=black,line width=0.8pt] (\x,0.9) rectangle (\x+0.7,0.9+\h);
  }
  \node[font=\tiny,text=white,rotate=90] at (10.7,3) {data data data data data};
  \foreach \x/\y in {1.2/1.5,1.6/4.2,10.2/5.6,0.8/3}{\node[star,star points=5,fill=yellow,draw=red,minimum size=0.6cm] at (\x,\y) {};}
  \node[font=\tiny,text=lime!50!black] at (5.4,0.45) {look at it!!!  best ever!!!  wow  wow  more wow};
\end{tikzpicture}
\end{center}

\begin{slopbox}
\textbf{(a)} What is the poster's main point? Can you tell?\\[3pt]\rule{\linewidth}{0.5pt}\\[6pt]
\textbf{(b)} Name two things that are not simple, because they only distract.\\[3pt]\rule{\linewidth}{0.5pt}\\[6pt]
\textbf{(c)} Name two things that are missing, so that the bars cannot be read.\\[3pt]\rule{\linewidth}{0.5pt}\\[6pt]
\textbf{(d)} Which part is not relevant to measuring at all?\\[3pt]\rule{\linewidth}{0.5pt}
\end{slopbox}

\subhead{Now you try}
On a separate sheet, redraw the poster with the same eight bars. Give it a title that says what was measured, label the axes with units, pick one or two calm colours, and delete everything that does not help. Check your poster against the three words: salient, simple, relevant.

\begin{center}
\begin{tikzpicture}
  \draw[ink!70,line width=1pt,rounded corners=3pt] (0,0) rectangle (11.4,1.0);
  \node[font=\sffamily\small,anchor=west] at (0.2,0.5) {Title:};
  \draw[line width=1.3pt] (1,1.6) -- (1,9.6);\draw[line width=1.3pt] (1,1.6) -- (11,1.6);
  \foreach \y in {2.6,3.6,...,8.6}{\draw[gray!40] (1,\y) -- (11,\y);}
  \node[font=\sffamily\scriptsize,rotate=90] at (0.3,5.6) {label and unit here};
  \node[font=\sffamily\scriptsize] at (6,1.3) {label what the bars stand for};
\end{tikzpicture}
\end{center}

\clearpage
\thispagestyle{fancy}
\subhead{Mixed review}
Each puzzle uses something from an earlier chapter.

\begin{puzzle}{Too many digits}
A student writes: ``My pencil is 17.43928~cm long.'' The ruler is marked in millimetres.\\
What is wrong, and what should the student write? \wline[3cm]\\
\rule{\linewidth}{0.5pt}
\end{puzzle}

\begin{puzzle}{Round, average, report}
A string was measured five times, in centimetres: 24.6, 24.8, 24.7, 24.6, 24.9.\\
(a) Total: \wline[1.6cm] \quad (b) Average: \wline[1.6cm] \quad (c) Average to the nearest tenth: \wline[1.6cm]\\
(d) Spread: \wline[1.6cm]\\
(e) Write the answer as one honest sentence.\\
\rule{\linewidth}{0.5pt}
\end{puzzle}

\begin{puzzle}{Match the problem to the fix}
Draw a line from each problem to the best fix.
\begin{center}
\begin{tikzpicture}[font=\sffamily\small]
  \foreach \y/\t in {0/{(a) Each reading is a little different},-1.1/{(b) A scale shows 4 g with nothing on it},-2.2/{(c) The mark seems to move as you lean},-3.3/{(d) The calculator shows nine digits}}{
    \node[draw=ink,rounded corners=2pt,text width=5.4cm,inner sep=4pt] at (0,\y) {\t};
  }
  \foreach \y/\t in {0/{look straight on},-1.1/{measure many times, average},-2.2/{round to what the tool can show},-3.3/{subtract the zero error}}{
    \node[draw=ink,rounded corners=2pt,text width=5.0cm,inner sep=4pt] at (7,\y) {\t};
  }
\end{tikzpicture}
\end{center}
\end{puzzle}

\clearpage
\thispagestyle{fancy}
\subhead{My measurement log}
Write down things you measure. Always include the unit and the tool.
\renewcommand{\arraystretch}{2.1}
\begin{center}\sffamily\scriptsize
\begin{tabularx}{\textwidth}{|p{1.5cm}|X|p{2.2cm}|p{1.8cm}|X|}
\hline
\textbf{Day} & \textbf{What I measured} & \textbf{Tool} & \textbf{Reading and unit} & \textbf{Anything odd?}\\ \hline
 & & & & \\ \hline
 & & & & \\ \hline
 & & & & \\ \hline
 & & & & \\ \hline
 & & & & \\ \hline
 & & & & \\ \hline
 & & & & \\ \hline
 & & & & \\ \hline
 & & & & \\ \hline
 & & & & \\ \hline
 & & & & \\ \hline
 & & & & \\ \hline
\end{tabularx}
\end{center}

\clearpage
\thispagestyle{fancy}
\subhead{Your turn: write a five-line report}
Choose any measurement from this book: the desk, the beans, the falling paper, the catch, the pencil. Fill in the five lines so a stranger could repeat what you did.
\begin{tcolorbox}[enhanced,colback=white,colframe=sky,boxrule=1.2pt,arc=2mm,left=8pt,right=8pt,top=6pt,bottom=6pt,before skip=7pt,after skip=7pt]
\sffamily\small
\textbf{1. What I measured.}\\[14pt]\rule{\linewidth}{0.5pt}\\[10pt]
\textbf{2. How I measured it.}\\[14pt]\rule{\linewidth}{0.5pt}\\[14pt]\rule{\linewidth}{0.5pt}\\[10pt]
\textbf{3. My readings (with units).}\\[14pt]\rule{\linewidth}{0.5pt}\\[10pt]
\textbf{4. My answer, rounded sensibly, with the spread.}\\[14pt]\rule{\linewidth}{0.5pt}\\[10pt]
\textbf{5. What I kept and what it hides.}\\[14pt]\rule{\linewidth}{0.5pt}\\[14pt]\rule{\linewidth}{0.5pt}\\[4pt]
\end{tcolorbox}
Check your report against the trust box questions. What did you check? What does that give you reason to trust? What does it not show?

\clearpage
\thispagestyle{fancy}
\subhead{What did we keep, across the whole book?}
For each chapter, say one thing that was simplified, what was kept, and what was set aside. A sample is done for you.

\renewcommand{\arraystretch}{2.3}
\begin{center}\sffamily\scriptsize
\begin{tabularx}{\textwidth}{|>{\bfseries}p{2.1cm}|X|X|X|}
\hline
\textbf{Chapter} & \textbf{What we simplified} & \textbf{What we kept} & \textbf{What we set aside}\\ \hline
1 Two answers & a desk length written as one number & the number, unit and how it was taken & small gaps, overlaps, judgements \\ \hline
2 Round on purpose & & & \\ \hline
3 Guess first & & & \\ \hline
4 Ten readings & & & \\ \hline
5 The instrument & & & \\ \hline
6 Reporting & & & \\ \hline
\end{tabularx}
\end{center}

\begin{multix}[wow]\multixbot[wow]\tcblower
Everything you have measured in this book was a simplification of something. Now you know how to say what was kept. Next time somebody gives you a number, you will want to ask: how was it taken?
\end{multix}
\clearpage
\thispagestyle{fancy}
\subhead{Questions I still have}
Write a question about measuring that this book made you wonder about. Then write how you might find out.
\par\vspace{1em}
\foreach \i in {1,...,12}{\noindent\rule{\linewidth}{0.4pt}\par\vspace{20pt}}
\clearpage
